\Needspace{8\baselineskip}\originalchapter{2006 作业二: 镜场、地磁与反转层}
题源: 19 Sep 2006, due 26 Sep.
\begin{exercise}\label{ex:06-2-1}
原题1. 轴对称镜 $B_z=B_0(1+z^2/A^2)$, $B_r=-B_0rz/A^2$. (a)中心沿z轴, $z=0$时能量W0=W\ensuremath{\perp}0+W\ensuremath{\parallel}0, 求z(t); (b)中心位于$z=0,r=A/2$时gradB与曲率漂移; (c)中心穿该点、能量同前, 求转向zt方程(不解), 需场线方程.
\end{exercise}
\begin{solution}
(a)\ensuremath{\mu}=W\ensuremath{\perp}0/B0, $m\ddot z=-2W_{\perp0}z/A^2$, 若$t=0$从$z=0$向$+z$, $z=A\sqrt{W_{\parallel0}/W_{\perp0}}\sin[\sqrt{2W_{\perp0}/(mA^2)}t]$. (b)该点\ensuremath{\nabla}|B|=0, gradB漂移零; 由场线 $\dd r/\dd z=-rz/(A^2+z^2)$ 得\ensuremath{\kappa}r=\ensuremath{-}r/A\ensuremath{{}^{2}}, 因而 $\mathbf v_c=-mv_\parallel^2r/(qB_0A^2)\mathbf e_\theta$. (c)场线 $r(z)=(A/2)A/\sqrt{A^2+z^2}$, 转向条件
\begin{equation}
 \sqrt{(1+z_t^2/A^2)^2+[r(z_t)z_t/A^2]^2}
 =\frac{W_0}{W_{\perp0}}.
\end{equation}
不能将off-axis轨道直接使用轴线上Bz替代总|B|.
\end{solution}
\begin{exercise}\label{ex:06-2-2}
原题2. 忽略太阳风影响, 地球偶极场 $\mathbf B=3\times10^{-5}(R_e/r)^3[-2\sin\theta\mathbf e_r+\cos\theta\mathbf e_\theta]$T, \ensuremath{\theta}纬度, Re\ensuremath{\simeq}6400km. (i)解dr/(r d\ensuremath{\theta})=Br/B\ensuremath{\theta}求场线; (ii)以赤道Be表|B|(\ensuremath{\theta}); (iii)赤道初W\ensuremath{\perp}=2W0/3,W\ensuremath{\parallel}=W0/3, 用最低非零小\ensuremath{\theta}项估镜纬度; (iv)以赤道漂移代表整段, 求环地周期, 1MeV质子在2Re实算并检查rL条件.
\end{exercise}
\begin{solution}
(i) $r=r_{eq}\cos^2\theta$. (ii) $B/B_e=\sqrt{1+3\sin^2\theta}/\cos^6\theta$. (iii)展开为1+9\ensuremath{\theta}\ensuremath{{}^{2}}/2, 转向B/Be=3/2得\ensuremath{\theta}\ensuremath{\simeq}1/3rad\ensuremath{\simeq}19.1°; 不是精确方程的数值根. (iv)赤道\ensuremath{\nabla}B=\ensuremath{-}3B/r er、曲率大小3/r, 漂移模 $v_D=(3W_\perp+6W_\parallel)/(|q|Br)=4W_0/(|q|Br)$. 两项同向, 对不同电荷方向反转. 周期 $2\pi r/v_D=965$s\ensuremath{\simeq}16.1min, Be=$3.75\times10^{-6}$T. $r_L=3.15\times10^4$m, 远小于r=$1.28\times10^7$m及曲率半径r/3, 可用漂移近似. 实际bounce平均与真实地磁会修正周期.
\end{solution}
\begin{exercise}\label{ex:06-2-3}
原题3. 剪切层题印为 $B_z=B_0$在y<\ensuremath{-}a, $B_z=B_0y/a$在|y|<a, $B_z=-B_0$在y>a. (i)中心初$x=0,y=a/2$,rL\ensuremath{\ll}a, 描述运动与W\ensuremath{\perp}限制; (ii)粒子初$x=y=0,v_y=V,v_x=0$, 写运动方程; (iii)最大|y|, 假定不到外均匀区, 给条件.
\end{exercise}
\begin{solution}
原件中层符号与外侧两端不连续. 本题轨道限于内层, 所以采用\textbf{印出的}内层 $B_z=by$, b=B0/a; 若改为连续的\ensuremath{-}B0y/a, 下列x方向反转而幅度条件不变. (i) $v_{\nabla B,x}=-\mu/(qy)$, 在a/2匀速沿x漂移, y保持; $r_L=2\sqrt{2mW_\perp}/(|q|B_0)\ll a$, 即 $W_\perp\ll q^2B_0^2a^2/(8m)$. (ii) $\dot v_x=(qb/m)yv_y$, $\dot v_y=-(qb/m)yv_x$, $\dot y=v_y$. (iii)积分 $v_x=qby^2/(2m)$, 能量 $v_x^2+v_y^2=V^2$, 故 $|y|_{\max}=\sqrt{2m|V|/(|q|b)}$. 不到外层要求 $m|V|<|q|B_0a/2$. 在零场不能用发散的1/B漂移代替这组精确内层轨道方程.
\end{solution}
